\section{AI 云调度难题：拓扑、数据与抽象层冲突}

前面指出通用云抽象与 AI workload 存在结构差异，本节把冲突落到调度和数据面。一个训练任务不仅需要若干 GPU，还需要这些 GPU 处于合适互联域、及时获得同一数据版本，并能在长作业失败后恢复；资源数量因此不能脱离位置和状态讨论。

\subsection{局部性约束重新变成一等公民}
传统 VM 语义尽量隐藏物理位置，分布式训练却会让跨交换机、跨机架或跨区域通信直接进入 step time。本小节定义 locality 的实际对象，包括 GPU 伙伴、数据缓存和网络路径，并说明拓扑感知调度为什么是性能与可靠性的共同要求。
在 1:09 左右，讲者指出 AI 训练任务对物理邻近性高度敏感。请求四台机器时，最好位于同机架并挂在同一交换机上，否则通信代价会吞噬算力收益。

\begin{importantbox}{旧云抽象与新负载矛盾}
传统云强调资源可替换与位置无关；AI 训练强调位置相关和链路稳定。两者目标不一致，导致通用抽象在 AI 场景下出现性能与成本折损。
\end{importantbox}

\subsection{数据面复杂度：不只是 ``装个框架''}
讲者提到，安装 PyTorch 很快，但把 PB 级数据接入目标存储与训练作业才是长期瓶颈。不同环境（对象存储、块存储、并行文件系统）在吞吐、延迟与一致性上差异巨大，数据管线设计成为系统成败关键。

\begin{knowledgebox}{数据面工程的三个核心动作}
\begin{enumerate}
    \item 数据布局：冷热分层、分片策略、预取窗口；
    \item 数据通路：对象存储到训练节点的高效拉取与缓存；
    \item 数据治理：版本、血缘、质量校验与回放能力。
\end{enumerate}
\end{knowledgebox}

\begin{table}[H]
\centering
\begin{tabular}{p{2.8cm}p{3.8cm}p{3.7cm}p{2.7cm}}
\toprule
\textbf{层面} & \textbf{传统云默认假设} & \textbf{AI 训练现实} & \textbf{改造方向} \\
\midrule
调度 & 位置无关，可替换 VM & 位置相关，拓扑敏感 & 拓扑感知调度器 \\
存储 & 统一接口即可 & 多存储栈性能差异显著 & 分层数据管线 \\
网络 & 服务间 RPC 为主 & AllReduce / 参数同步密集 & 高带宽低抖动网络 \\
运维 & 微服务弹性扩缩容 & 训练作业长时、失败代价高 & 作业级容错与恢复 \\
\bottomrule
\end{tabular}
\caption{AI 负载下 ``抽象便利性'' 与 ``物理现实'' 的矛盾}
\end{table}

\begin{warningbox}{平台团队常见低估：忽略数据搬运成本}
很多团队在成本估算时只看 GPU 单价，忽略数据准备、跨区传输、存储读放大和作业重试成本，导致总成本明显高于预期。
\end{warningbox}

\subsection{本章小结}
AI 系统设计正在从 ``计算资源管理'' 扩展为 ``计算+数据+拓扑'' 协同优化问题。调度与数据面能力将决定平台竞争力。

